\documentclass{article}
\usepackage{assignment_preamble}

\title{Homework 2}
\author{Ravi Kini}
\date{November 27, 2024}

\begin{document}

\maketitle

\problem
In terms of location of activation, bilingual speakers with higher proficiency experience activation in more overlapping areas, while bilinguals with lower profificiency experience activation in areas that are more spread out. In terms of strength of activation, speakers with higher proficiency experience less intense activation in areas related to language control like the prefrontal cortex due to requiring less processing power to parse the language.

\clearpage

\problem
Slower lexical access and a smaller vocabulary size are cited as disadvantages of bilingualism.

\clearpage

\problem
One cited cognitive advantage of bilingualism is improved multitasking. However, the topic of cognitive advantage for bilinguals is controversial because of the number of external factors like socioeconomic status or education that can impact these advantages, and the influence of the ceiling effect.

\clearpage

\problem
When interacting with non-native speakers, native speakers will simplify their speech by speaking slower and using less complex sentences, sometimes even omitting elements in a matter that results in ungrammatical sentences. This is done to maximize common ground between native and non-native speakers to make communication easier.


\clearpage

\problem
A multilingual society is a society where multiple languages are spoken.

\clearpage

\problem
Two languages may have different status in a society due to factors like social domination of one group over another, among other macrosocial and microsocial factors. For example, in India, the dominance of Hindi and English in comparison to state languages can be explained by the dominance of the north and the effects of colonialism, respectively.

\clearpage

\problem
When a bilingual speaker exhibits coordinate unbalanced bilingualism, having a lower proficiency in the L2 than in the L1, this results in some overlap in the areas used for storage and processing, with the L2 being activated in more areas in the right hemisphere and the activation being more diffused. We can detect this by measuring the activation of neurons in areas of the brain related to language while a bilingual speaker uses their L1 versus when they use their L2.

\clearpage

\problem
More code switching is likely to occur when a bilingual speaker is experiencing a lot of emotion, due to different languages having different ways of expressing emotions, usually using the L1 to express stronger emotions, and usually for younger or less proficient speakers. Less code switching is likely to occur as a speaker ages and gains more proficiency, as the speaker becomes more detached from the language.

\clearpage

\problem
An example of a less prestigious variant influencing a more prestigious variant is the transfer of evidentiality into Andean Spanish from Quechua and Aymara as a result of the numbber of bilingual speakers using evidentiality in their L1 language. On the other hand, in India, speakers of lower social classes seek to imitate the speech of higher classes, resulting in an evolution in their speech patterns over time.


\clearpage

\problem
In multilingual societies, acculturation in the process of adapting to the values of a community. For exmaple, immigrants to the United States acculturate by developing greater English proficiency, as many Americans are monolingual, and adhering to American social norms.

\clearpage

\problem
A lack of motivation can result in a failure to acquire the L2 language in a multilingual society; for example, L1 English speakers in Quebec failed to progress past a certain level of proficiency in French. Another reason could be attitudes towards the speakers by the dominant social group. For example, L1 Finnish students in Sweden did worse in their L2 language than those in Australia as they were viewed negatively as a minority group.

\clearpage

\problem
Social motivation impacts the level of proficiency one can gain in a language. For example, English L1 speakesr in Quebec showed limits to how much French proficiency they gained, even after exposure in early childhood, due to being motivated to not learn a second language.

\clearpage

\problem
Language is used to create identities in multilingual communities. For example, the Spanish conquerors of Mexico differentiated themselves from the local Indian population by refusing to teach them Castilian Spanish despite gaining proficiency in the indigenous language.

\clearpage

\problem
Usage of a L2 language by biilnguals has been shown to to result in more rational judgements in comparison to the L1 due to the increased emotional distance bilinguals have towards their L2 language.

\end{document}